% ============================================================
%  SEKCJA 6: RYNEK MIODU UE I IMPLIKACJE REGULACYJNE
% ============================================================

\section{Rynek miodu UE i implikacje regulacyjne}

Europejski rynek miodu jest największym regionalnym rynkiem
importowym na~świecie, a~jego strukturalne cechy, chroniczny
deficyt produkcji własnej, wysoka zależność od~importu
z~krajów trzecich i~szerokie zróżnicowanie jakościowe
dostarczanych produktów, stanowią kontekst, w~którym
decyzje regulacyjne dotyczące kryteriów dyskwalifikacji
miodu mają bezpośrednie konsekwencje ekonomiczne, handlowe
i~konsumenckie. Niniejsza sekcja analizuje ten kontekst
rynkowy jako uzasadnienie dla~postulowanego w~artykule
podejścia precyzyjnego, opartego na~dojrzałości funkcjonalnej
i~aktywności wody, w~miejsce podejścia nadmiernie
uproszczonego, opartego na~liczbie drożdży lub~kryterium
morfologicznym.

%,,,,,,,,,,,,,,,,,,,,,,,,,,,,,,
\subsection{Struktura i dynamika importu miodu do UE}
%,,,,,,,,,,,,,,,,,,,,,,,,,,,,,,

Unia Europejska jest największym regionalnym importerem
miodu na~świecie, generując deficyt bilansu handlowego
w~sektorze pszczelarskim, który nie~ma odpowiednika
w~żadnym innym regionie~\cite{cbi2024,pollinatorhub2025}.
Produkcja własna 27~krajów członkowskich wynosi szacunkowo
272,286~tys.~ton rocznie (dane 2022,2023), co~pokrywa
jedynie 55,60\% krajowego zapotrzebowania, szacowanego
na~ok.~480,500~tys.~ton~\cite{pollinatorhub2025,ec_honey_autumn2024}.
Różnicę tę~wypełnia import z~krajów trzecich.

Według danych Eurostatu za~rok 2023, UE sprowadziła
\textbf{163~700~ton} miodu z~krajów spoza UE o~wartości
\textbf{359,3~mln~EUR}; eksport wyniósł 24~900~ton
(146~mln~EUR)~\cite{eurostat_honey2023}. Dodając
handel wewnątrzwspólnotowy (miód przemieszczany między
państwami członkowskimi po~uprzednim imporcie z~zewnątrz),
łączny wolumen miodu importowanego do~łańcucha dostaw UE
szacowany jest na~\textbf{300,360~tys.~ton} rocznie~\cite{cbi2024}.
W~2024~roku łączny import (extra-UE + wewnątrzwspólnotowy)
wyniósł ok.~290~tys.~ton (spadek o~3,1\% r/r po~trzyletnim
wzroście), lecz długoterminowy trend pozostaje rosnący:
rynek UE ma~osiągnąć wolumen 445~tys.~ton do~2035~roku
(CAGR +1,9\%)~\cite{eurostat_honey2023}.

Głównymi dostawcami zewnętrznymi w~2023~roku były:
\textbf{Chiny} (37\% importu extra-UE), \textbf{Ukraina} (28\%),
\textbf{Argentyna} (ok.~7\%), \textbf{Meksyk} i~\textbf{Kuba}~\cite{tridge2024}.
Największymi krajowymi importerami wewnątrz UE były Niemcy
(41~tys.~ton, 25\% importu extra-UE), Belgia (31,4~tys.~ton,
19\%), Polska (23,3~tys.~ton, 14\%), Hiszpania (15,7~tys.~ton)
i~Francja (7,7~tys.~ton)~\cite{eurostat_honey2023}. 

%,,,,,,,,,,,,,,,,,,,,,,,,,,,,,,
\subsection{Skala fałszowania miodu na rynku UE: dane empiryczne}
%,,,,,,,,,,,,,,,,,,,,,,,,,,,,,,

Skala problemu autentyczności miodu na~rynku UE jest
dokumentowana przez szereg badań prowadzonych na~poziomie
instytucjonalnym. Skoordynowana akcja kontrolna Komisji
Europejskiej ,,From the Hives'' (2021,2022), obejmująca
analizę 320~próbek miodu importowanego, wykazała, że~\textbf{46\%}
badanych próbek było podejrzanych o~niezgodność z~Dyrektywą
Miodową, co~oznacza trzykrotny wzrost w~porównaniu z~analogicznym
badaniem z~lat 2015,2017, kiedy wskaźnik ten wynosił~14\%~\cite{eufromhives2021}.
Najwyższy bezwzględny odsetek podejrzanych przesyłek
dotyczył Chin (74\%), natomiast Turcja wykazała najwyższy
względny wskaźnik (93\%)~\cite{eufromhives2021}. Wśród
zidentyfikowanych praktyk fałszerskich dominowało dodawanie
syropów cukrowych (z~trzciny cukrowej, kukurydzy, ryżu)
oraz blendowanie miodu autentycznego z~produktami
syropowymi w~proporcjach poniżej progu detekcji
metodami konwencjonalnymi~\cite{eufromhives2021,schoder2026}.

Dane te~wskazują na~realny i~poważny problem autentyczności
miodu na~rynku UE, lecz jednocześnie wyraźnie profilują
dominującą formę fałszowania jako \emph{adulterację syropami
cukrowymi}, nie~zaś jako dystrybucję miodu z~podwyższoną
liczbą drożdży. Dokumenty JRC i~OLAF nie~wskazują
na~żaden przypadek dyskwalifikacji miodu wyłącznie
na~podstawie liczby CFU/g drożdży bez~jednoczesnego
stwierdzenia aktywnej fermentacji lub~adulteracji
syropowej~\cite{eufromhives2021}. Jest to~istotne
spostrzeżenie z~perspektywy niniejszego artykułu:
praktyka fałszowania skupia się na~chemicznym rozcieńczaniu
miodu, nie~na~manipulacji jego mikroflorą.

W~odpowiedzi na~ten kryzys Komisja Europejska przyjęła
Dyrektywę (UE) 2024/1438 (część pakietu ,,Dyrektyw
śniadaniowych''), wprowadzającą od~2026~roku obowiązek
deklarowania wszystkich krajów pochodzenia składników
mieszanek miodowych w~głównym polu wizualnym
etykiety~\cite{cbi_traceability2024}. Wdrożenie tego
wymogu wzmacnia infrastrukturę traceability, lecz
nie~rozwiązuje problemu braku walidowanych parametrów
specyficznych dla~niedojrzałości.

%,,,,,,,,,,,,,,,,,,,,,,,,,,,,,,
\subsection{Luki analityczne: brak walidowanego biomarkera niedojrzałości}
%,,,,,,,,,,,,,,,,,,,,,,,,,,,,,,

Identyfikacja miodu niedojrzałego jako odrębnej kategorii
fałszowania wymaga istnienia narzędzi analitycznych zdolnych
do~wiarygodnego rozróżnienia między miodem niedojrzałym
a~dojrzałym, i~tutaj nauka napotyka na~fundamentalne
ograniczenie. Jak wykazał przegląd Tsagkaris i~in.~\cite{tsagkaris2021},
standardowe metody fizykochemiczne (wilgotność, diastaza,
inwertaza, HMF, profil cukrowy) nie~są wystarczające
do~jednoznacznego orzeczenia o~niedojrzałości, gdyż:

\begin{enumerate}
    \item \textbf{Brak jest markera specyficznego:}
          żaden pojedynczy związek chemiczny ani~enzym
          nie~jest identyfikowany jako biomarker wyłącznie
          niedojrzałości w~sposób niezależny od~typu
          botanicznego i~geograficznego; parametry
          fizyczno-chemiczne dojrzałego miodu z~tropiku
          mogą przypominać parametry miodu niedojrzałego
          ze~strefy umiarkowanej~\cite{tsagkaris2021,guo2020};
    \item \textbf{Metagenomika i~metabolomika mają ograniczoną
          moc różnicującą bez~baz referencyjnych:}
          jak podkreśla Schoder~\cite{schoder2026}
          w~przeglądzie z~2026~roku, wysoka rozdzielczość
          analityczna metod omicznych (genomika, proteomika,
          NMR, GC-MS, LC-HRMS) nie~gwarantuje jednoznacznej
          atrybucji fałszowania bez~solidnych, globalnie
          reprezentatywnych baz referencyjnych; \textit{,,high
          analytical resolution alone therefore does not
          guarantee unambiguous fraud attribution''}~\cite{schoder2026};
    \item \textbf{Efekty blendu maskują niedojrzałość:}
          blendowanie miodów niedojrzałych z~dojrzałymi
          w~proporcjach poniżej 50\% jest analitycznie
          nieodróżnialne od~miodu dojrzałego metodami
          konwencjonalnymi~\cite{schoder2026},należy podkreślić, iż problem ten dotyczy również etapu pozyskiwania miodu z plastrów, nie moża określić jaka część miodu na czas wirowani była dojrzała a jaka niedojrzała;
    \item \textbf{Obecność drożdży nie~jest markerem
          specyficznym:} jak wykazano w~sekcji~4,
          liczba drożdży osmofilnych w~miodzie jest
          determinowana głównie przez warunki klimatyczne
          produkcji i~systemy ula, a~nie przez stopień
          dojrzałości, co~wyklucza CFU/g jako
          użyteczny biomarker niedojrzałości~\cite{rodriguezmachado2024,schoder2026}.
\end{enumerate}

Najnowsze badania wskazują na~obiecujące kierunki
w~poszukiwaniu biomarkerów niedojrzałości: Sun i~in.~\cite{sun2021}
zidentyfikowali kwas dekenodwukarboksylowy i~inne kwasy
tłuszczowe pszczelego pochodzenia jako potencjalne markery
molekularne, których stężenie koreluje ze~stopniem dojrzałości
w~sposób niezależny od~liczby drożdży. Metabolomika
oparta na~NMR~\cite{nmrhoney2023} umożliwia jednoczesną
kwantyfikację kilkudziesięciu związków w~jednym pomiarze,
co~daje podstawy do~opracowania wielowymiarowego wskaźnika
dojrzałości. Jednak do~dnia dzisiejszego żaden z~proponowanych
biomarkerów nie~przeszedł pełnej walidacji wielolaboratoryjnej
z~uwzględnieniem globalnej różnorodności typów botanicznych
i~geograficznych~\cite{tsagkaris2021}.

Trzeba mieć na uwadze, że prawnie ustanowione metody winny być możliwie proste, szybkie i dostępne gdyż będą stanowić podstawowe narzędzie każdego pszczelarza (podobnie jak obecnie refraktometr) do określenia terminu rozpoczęcia wirowania gdyż to ten etap jest decyduje czy miód jest dojrzały czy dojrzały nie jest.  

%,,,,,,,,,,,,,,,,,,,,,,,,,,,,,,
\subsection{Ryzyko regulacyjne: konsekwencje nadmiernie rygorystycznej interpretacji}
%,,,,,,,,,,,,,,,,,,,,,,,,,,,,,,

Przyjęcie kryterium dyskwalifikacyjnego opartego na~liczbie
drożdży CFU/g, bez~uwzględnienia aktywności wody,
warunków klimatycznych produkcji i~braku aktywnej fermentacji
, niesie ze~sobą szereg poważnych ryzyk regulacyjnych
i~rynkowych, których analiza jest konieczna dla~właściwej
oceny implikacji stanowiska Apimondia:

\paragraph{Ryzyko dyskryminacji producentów tropikalnych.}
Jak wykazano w~sekcji~5, miody z~systemów tropikalnych
i~tradycyjnych naturalnie zawierają wyższe stężenia
drożdży osmofilnych ze~względu na~typ ula, warunki
klimatyczne i~bogatszą endemiczną bioróżnorodność
mikrobiologiczną~\cite{tadele2025,stinglessyeastbrazil2021}.
Dyskwalifikacja tych miodów na~podstawie CFU/g
skutkowałaby de~facto wykluczeniem z~rynku UE produktów
legalnie produkowanych w~krajach Afryki Wschodniej,
Azji Południowo-Wschodniej i~Ameryki Łacińskiej,
jednocześnie nie~dotykając miodów adulterowanych
syropami cukrowymi (które zawierają standardową
lub~obniżoną liczbę drożdży po~procesie
ultrafiltracji)~\cite{eufromhives2021}.

\paragraph{Ryzyko pogłębienia deficytu i wzrostu cen.}
Wykluczenie z~rynku UE miodów z~Chin, Ukrainy i~Argentyny
, trzech największych dostawców odpowiadających łącznie
za~$\approx 72\%$ importu extra-UE~\cite{tridge2024}
, w~oparciu o~kryterium CFU/g spowodowałoby gwałtowny
wzrost cen miodu na~rynku europejskim i~pogłębienie
deficytu, szczególnie dotkliwe dla~konsumentów w~krajach
o~niskim spożyciu własnym (Polska, Niemcy, Belgia
jako netto importerzy)~\cite{eurostat_honey2023}.

\paragraph{Ryzyko tworzenia arbitralnych barier handlowych.}
Stosowanie kryterium niemającego podstawy w~Codex Alimentarius
ani~w~Dyrektywie UE 2001/110/WE może być kwestionowane
przez partnerów handlowych w~ramach Porozumienia WTO-SPS
(Umowa o~stosowaniu środków sanitarnych i~fitosanitarnych),
które wymaga, aby środki ograniczające handel miały
podstawę naukową i~były proporcjonalne do~zidentyfikowanego
ryzyka~\cite{codex_stan12,eudir2001110}.
Ryzyko pozyskiwania miodu niedojrzałego występuje we wszystkich miodach, niezależnie od miejsca prowadzenia gospodarki pasiecznej. Wprowadzone wymagań i ograniczenia winny mieć wymiar globalny i swoje proporcjonalne  odzwierciedlenie w prawie unijnym wobec unijnych producentów miodu. 


\paragraph{Ryzyko rozminięcia z rzeczywistym problemem.}
Jak wykazała akcja ,,From the Hives'', dominującą
formą fałszowania miodu w~UE jest adulteracja
syropami cukrowymi, a~nie dystrybucja miodu
niedojrzałego o~podwyższonej liczbie drożdży~\cite{eufromhives2021}.
Skupienie zasobów kontrolnych na~pomiarze CFU/g
kosztem analizy izotopowej ($\delta^{13}$C,
$\delta^{2}$H, $\delta^{18}$O) i~NMR fingerprinting
prowadziłoby do~alokacji zasobów inspektorskich
niezgodnie z~rzeczywistym profilem ryzyka fałszowania.
Warto jednak zaznaczyć, że Komisja Europejska później wycofała się z traktowania profilowania H-NMR jako wystarczająco ugruntowanej metody urzędowej kontroli, podkreślając potrzebę dalszych badań, walidacji i~harmonizacji metod analitycznych\cite{EuropeanCommission2023NMR}.


%,,,,,,,,,,,,,,,,,,,,,,,,,,,,,,
\subsection{Postulaty regulacyjne: harmonizacja, traceability i proporcjonalność}
%,,,,,,,,,,,,,,,,,,,,,,,,,,,,,,

Na~podstawie powyższej analizy rynkowej i~regulacyjnej
niniejszy artykuł formułuje następujące postulaty
adresowane do~organów normalizacyjnych (ISO, Codex),
instytucji UE (DG SANTE, KE) i~organizacji branżowych
(Apimondia, Copa-Cogeca):

\begin{enumerate}
    \item \textbf{Aktualizacja Dyrektywy UE 2001/110/WE
          o~kryterium aktywności wody ($a_w$):}
          wprowadzenie obowiązkowego pomiaru $a_w$
          obok wilgotności jako parametru pierwszorzędowego
          oceny ryzyka fermentacji. Pomiar $a_w$ jest
          metodą zwalidowaną, dostępną komercyjnie
          (przyrządy typu AquaLab, Rotronic HygroLab)
          i~niedrogą; stanowiłoby to~substytucję
          kryterium morfologicznego na~kryterium
          fizykochemiczne o~bezpośrednim uzasadnieniu
          naukowym~\cite{ruegg1981,zamora2006};

    \item \textbf{Ustanowienie wieloparametrycznego protokołu
          oceny mikrobiologicznej:} ocena ryzyka
          fermentacji powinna być prowadzona
          równocześnie dla~$a_w$, zawartości wody,
          temperatury przechowywania i~liczby drożdży
          (zgodnie z~równaniem ryzyka (\ref{eq:fermentation_risk})
          zaproponowanym w~sekcji~4.3); żaden z~tych
          parametrów nie~może być rozstrzygający
          w~izolacji od innych~\cite{analytica2020,apinz2021};

%     \item \textbf{Wdrożyć i~obowiązkowo stosować
%           systemy traceability w~całym łańcuchu dostaw:}
%           Dyrektywa (UE) 2024/1438 stanowi krok
%           w~dobrym kierunku, lecz niewystarczający;
%           wymagana jest identyfikowalność do~poziomu
%           pasieki z~dokumentacją wilgotności i~temperatury
%           na~każdym etapie \textit{post-harvest}
%           (cold chain monitoring)~\cite{cbi_traceability2024,tsagkaris2021};

%     \item \textbf{Budować globalne bazy referencyjne
%           autentyczności miodu:} jak wskazuje PMC~2026~\cite{schoder2026},
%           zaawansowane metody omiczne są nieefektywne
%           bez~reprezentatywnych baz referencyjnych;
%           projekt EU Reference Laboratory for Honey
%           (EURL) powinien zostać rozszerzony o~próbki
%           z~systemów tropikalnych i~tradycyjnych;

%     \item \textbf{Rozróżniać regulacyjnie trzy kategorie
%           produktów:} (a)~miód dojrzały funkcjonalnie
%           z~naturalną mikroflorą drożdżową, wolny
%           od~ograniczeń; (b)~miód z~zastosowaną
%           technologiczną dehumidyfikacją, dopuszczony
%           z~deklaracją procesu; (c)~miód systemowo
%           niedojrzały lub~adulterowany, zakazany;
%           wyłącznie ta~kategoryzacja jest proporcjonalna
%           do~rzeczywistego profilu ryzyka
%           na~rynku~\cite{garcia2025,tadele2025,codex_stan12}.
\end{enumerate}